\chapter{Literature Review}
\label{ch:literature-review}

\section{Multimodal mock-interview scoring}

Naim et al. introduced automated analysis of mock job interviews using verbal, vocal and visual
behaviour~\cite{Naim2015}, then reported an expanded treatment of interview-performance
prediction~\cite{Naim2018}. That work motivates the present corpus and feature boundary. It does
not turn a crowd rating into an observed employment outcome. Related studies compare asynchronous
and face-to-face communication-skill measurement~\cite{Rasipuram2016}. Automated video-interview
personality assessments have separately been tested for reliability, validity and
generalisation~\cite{Hickman2022}. Together, this literature makes construct validity central:
prediction of a rater aggregate is evidence about agreement with that aggregate, not validation of
hiring decisions.

Employment-interview validity is conditional on what is asked, how responses are scored and
which later criterion is tested. McDaniel et al.'s meta-analysis established a foundation for
examining interview validity~\cite{McDaniel1994InterviewValidity}. Levashina et al. review the
design and evidence for structured interviews~\cite{Levashina2014Structured}, while Wingate et
al. estimate validity for distinct constructs rather than treating every interview rating as the
same measure~\cite{Wingate2025ConstructValidity}.
Re-estimates of personnel-selection validity also caution against carrying inflated summary
coefficients into a new setting~\cite{Sackett2022Validity}. The supplied MIT label is an
aggregate judgement about a mock interview, with no observed job-performance criterion.
Consequently, even a perfectly reproduced predictive score would leave the employment-validity
question open.

The target is also an operational measure inside a social process, not the construct itself.
Measurement theory distinguishes an abstract concept from its observed proxy
~\cite{JacobsWallach2021}. Selbst et al. show how abstraction from an actual institution can
conceal social context~\cite{Selbst2019}. Suresh and Guttag trace possible sources of harm across
the machine-learning life cycle~\cite{SureshGuttag2021}. Studies of commercial hiring tools
caution against accepting broad fairness claims without checking what is predicted and how it is used
~\cite{Raghavan2020}. An external audit of hiring-related personality predictions examined
stability~\cite{Rhea2022}, while an evaluation of algorithmic recruitment-video analysis found
that accuracy and group disparity could coexist~\cite{Kochling2021}. Both concern different models,
data and targets. They motivate the present audit but cannot make its crowd rating an employment
outcome or turn its observed score gaps into evidence of deployed discrimination.

Visual and vocal cues cannot simply be dismissed as noise or assumed to measure competence.
DeGroot and Motowidlo examined their influence on interviewers' judgements and relation to job
performance~\cite{DeGroot1999Cues}; a later meta-analysis synthesised evidence on nonverbal
cues across interviews~\cite{MartinRaugh2023Nonverbal}. Human judgments can also contain implicit
bias~\cite{Purkiss2006InterviewBias}, and criterion-related validity has varied among interviewers
and panels~\cite{VanIddekinge2006Interviewers}. Those findings support a
two-sided interpretation here: modality signal may help predict the recorded rater aggregate,
while neither signal use nor rater agreement proves a fair, job-relevant construct. The MIT
target audit assesses agreement among its outcome raters, but cannot establish which cues were
valid predictors of later employment outcomes.

The source of a rating is part of the measurement problem. In a study of video-interview
personality assessments, Hickman et al. found that evidence of reliability and validity differed
between models trained against interviewer reports and those trained against self-reports
~\cite{Hickman2022}. That finding is not a performance benchmark for this study: personality and
crowd-perceived hireability are different targets. Booth et al. likewise analyse multimodal
hireability inference through a psychometric account of bias, drawing attention to the distinction
between the behaviour recorded, the judgement elicited from a rater and the model that learns that
judgement~\cite{Booth2021}. Consequently, a small error against the supplied rating would not, on
its own, establish that the rating measures job-relevant ability or that a model should inform
selection.

Recent automated-assessment work has tested psychometric properties of specific video-interview
competencies~\cite{Liff2024Competency} and the reliability, validity and bias of automated
cognitive-ability assessments~\cite{Hickman2025Cognitive}. Another study examines the transfer
of personality scores from low-stakes verbal data to high-stakes interviews
~\cite{Stevenor2024HighStakes}. Booth et al. explicitly combine psychometric and computational fairness perspectives on
automated video interviews~\cite{Booth2021Psychometrics}. This literature offers a standard
for what stronger validation would require, not a surrogate validation dataset for the present
crowd-perceived score. In particular, transfer from a low-stakes verbal task to a high-stakes
interview is an empirical question rather than an automatic property of a trained encoder.

Interview format also changes the social and measurement setting. Research on digital interviews
examines applicant reactions and interviewer ratings~\cite{Langer2017DigitalInterviews}.
Videoconference-interview research examines factors contributing to differences from
face-to-face ratings and perceptions~\cite{Basch2021Camera}. For asynchronous interviews,
Brenner et al. study the applicant's perspective~\cite{Brenner2016AVIReactions}; Lukacik et al.
set out design and use questions~\cite{Lukacik2022AVIDesign}; and Suen et al. test how synchrony
and AI framing affect ratings and attitudes~\cite{Suen2019SynchronyAI}.
Applicant reactions to algorithmic rather than recruiter evaluation have also been studied
~\cite{Oostrom2024ApplicantReactions}. These are deployment and format-validity concerns;
the MIT mock-interview data cannot estimate those reactions or isolate a format effect.

Gucluturk et al. model first impressions from short public videos with multimodal residual
networks~\cite{Gucluturk2018}. Escalante et al. examine modelling, recognition and explanation
of apparent personality from videos~\cite{Escalante2022}. They are discussed here only as
adjacent literature. Their benchmark is not the study corpus, is not included in the analysis and
does not supply targets or protected annotations for this thesis.

The MIT material has also supported a distinct multimodal feedback task. Agrawal et al. used its
69 candidates and 138 pre/post interviews to classify behavioural ratings such as engagement,
speaking rate and eye contact~\cite{Agrawal2020}. This is a closer data precedent than a generic
video benchmark, but its classified cues are not this thesis's continuous
\emph{RecommendHiring} outcome and its reported accuracy is not a directly comparable baseline.
Other automated-interview systems have been evaluated on separately collected online interviews
~\cite{Chen2017}, with window-consistency multimodal scoring for asynchronous video
~\cite{Lv2024} and psychometric analysis of asynchronous interview ratings
~\cite{Koutsoumpis2024}. Their differing interview formats, targets and populations motivate
comparison of measurement choices rather than numerical
ranking against the present 69-person corpus.

Earlier computational studies illustrate why the interview setting is not a single prediction
task. Nguyen et al. related nonverbal behaviour in employment interviews to evaluators'
hirability judgements~\cite{Nguyen2014HireMe}; Chen et al. scored one-way monologues
~\cite{Chen2016Monologue}; and HireNet modelled questions in asynchronous applications
~\cite{Hemamou2019HireNet}.
These studies differ in who judges the interview, which behaviours are available and whether the
target is a score or a recruiter decision. Their model accuracies cannot be imported as evidence
that the present crowd rating measures job performance. A closer recent MIT comparison uses
pretrained transformers and facial embeddings alongside a second interview corpus
~\cite{Fithian2025}; its cross-context results help frame modality dependence, but do not
replicate this thesis's frozen-feature, proxy-group or audited-disparity design. Transfer from
interviews to oral-presentation assessment has likewise been examined in a related small-data
setting~\cite{Tun2023Transfer}, not established for this thesis's outcome.

Recording conditions are an additional measurement pathway. Controlled experiments with video
interviews found that audiovisual quality changed assessments of candidates
~\cite{Fiechter2018}. That finding makes lighting, framing and sound quality plausible
nuisance factors in visual and audio representations; it does not establish that they caused
any observed gap in the MIT recordings. The appropriate response here is to state the input
boundary and avoid interpreting a representation cluster as an applicant attribute.

\section{Multimodal representation learning}

Modern frozen encoders permit a small study to reuse large-scale representation learning without
fine-tuning millions of parameters. DINOv2 produces general-purpose visual representations
~\cite{Oquab2024}; WavLM learns speech representations from self-supervised pretraining
~\cite{Chen2022WavLM}; MPNet combines masked and permuted pre-training
~\cite{Song2020MPNet}; Sentence Transformers describes the sentence-embedding framework
~\cite{Reimers2019}, and the exact contrastively tuned MPNet-based checkpoint used here has its
own model record~\cite{SentenceTransformersAllMPNet2025}. Their scale does
not eliminate dataset limitations. Encoder lineage, input selection and pooling can encode
appearance, active-interval mixed-speaker acoustic and linguistic correlates unrelated to the intended score
construct.

The choice of frozen features and a small fitted head has methodological precedent, not a
guarantee of superiority. Kornblith et al. examine transfer of image-model representations
~\cite{Kornblith2019Transfer}; Tian et al. study simple baselines in few-shot classification
~\cite{Tian2020FewShot}. Both make the target domain and evaluation design important. Those
image-classification results do not measure DINOv2 performance on candidate-face crops or
justify transferring their accuracy estimates to continuous interview scoring.

The present design therefore freezes all three encoders, stores their exact revisions and uses a
small, shared downstream backbone. Fresh representations are the primary inputs. Corpus-supplied
facial, SHORE and prosodic features are retained for coverage checks and secondary explanations,
not silently mixed into the primary model.

Fusion architectures in neighbouring multimodal work include tensor interactions
~\cite{Zadeh2017TFN} and gated units~\cite{Arevalo2017GMU}. A later survey also covers
attention- and sequence-based approaches~\cite{Zhao2024}. Their capacity and temporal alignment
assumptions are not interchangeable with this study's
low-capacity, fixed-fusion scorer. Changing the fusion mechanism would also change capacity and
alignment in the proxy-versus-sampling comparison on only 69 independent participants. The fixed
backbone is an experimental control, not a claim that simple fusion is generally optimal.

The three inputs also have different evidential boundaries. A candidate-face crop is an appearance
measurement, the candidate-only transcript is a selected record of speech content, and the
active-interview audio includes both candidate and interviewer speech. A frozen encoder cannot
separate these sources merely because its parameters are fixed. This matters when an apparent
modality effect is interpreted: an audio-derived cluster may reflect the interviewer or recording
conditions as well as the candidate. The present work therefore treats modality diagnostics as
evidence about these representations, not as proof of a participant trait.

Nor does freezing make an embedding socially neutral. Soft-biometric information has been
recovered from face-recognition embeddings~\cite{Terhorst2021SoftBiometric}. Gender and language
signals have been studied in pretrained speech representations~\cite{Guillaume2024Speech}.
Association tests have exposed biases or measurement problems in sentence embeddings
~\cite{May2019SentenceBias} and contextual word embeddings~\cite{Kurita2019ContextBias}.
These papers concern different face, speech and text models from DINOv2, WavLM and MPNet. They
motivate caution about possible encoded cues, not a claim that the exact frozen checkpoints
used here leaked any particular attribute at a measured rate.

\section{Fairness evidence with and without demographic labels}

Fairness-aware video-interview research has documented performance and score differences across
annotated groups~\cite{Booth2021}. Kim et al. evaluate mitigation with an observed sensitive
label~\cite{Kim2023}. Hemamou et al. use indirect facial signals in an adversarial interview
model~\cite{HemamouIndirect2021}; Putra et al. adapt adversarial reweighting to multimodal
interview assessment~\cite{Putra2024}. These studies establish relevant methods, but they do not
make a demographic-free intervention equivalent to a protected-group intervention.

There is a longer methodological lineage behind these comparisons. Fair representations seek to
retain task information while limiting association with a designated sensitive variable
~\cite{Zemel2013}. Gradient reversal was developed for domain-adversarial training
~\cite{Ganin2015}; adversarial censoring seeks to remove a designated signal from a
representation~\cite{Edwards2016}, and adversarial fairness objectives explicitly target
unwanted bias~\cite{ZhangBL2018}. These approaches clarify what the present study does
\emph{not} assert: a cluster fitted without a sensitive label is
not evidence that a sensitive variable has been removed from the representation. Prior work on
multimodal apparent-personality assessment has also evaluated bias-mitigation strategies on
different annotated video corpora~\cite{Yan2020}; its results remain adjacent evidence, not an
MIT-specific fairness result.

The distinction is visible in the objectives. Kim et al. optimise a Wasserstein score-distribution
penalty and an adversarial representation objective using an available binary sensitive-attribute
annotation~\cite{Kim2023}. Hemamou et al. instead use face representations as indirect adversarial
targets without collecting the sensitive attribute for training, then evaluate recoverability
against labelled groups~\cite{HemamouIndirect2021}. Neither procedure is the experiment here:
the present W1 penalty uses clusters fitted separately from visual or interview-audio features,
and the perceived-attribute labels are withheld until the ordinary predictions are fixed. The
comparison borrows a distributional objective, not another study's protected-group estimand,
architecture or empirical result.

For regression, fairness definitions and reduction-based algorithms specify
the quantity constrained and the sensitive attribute observed~\cite{Agarwal2019FairRegression}.
Zhao examines the costs and benefits of fair regression under an explicit fairness objective
~\cite{Zhao2023FairRegression}. Petersen et al. examine fairness assessment for risk-score
models~\cite{Petersen2023RiskScores}; Corbett-Davies et al. show more broadly how the choice of
measure can misstate fairness~\cite{CorbettDavies2023}. Thus distributional parity, calibration,
error balance and eventual decision rules are distinct questions.
This thesis reports a continuous predicted-score contrast; it has no decision threshold,
calibrated employment risk or observed job-performance criterion. Thus its W1 result is not a
complete assessment of a hiring decision, even if the audited categories were measured without
error.

ARL searches for high-loss regions without demographic annotations by learning adversarial
example weights~\cite{Lahoti2020}. It is an appropriate generic comparator because it does not
require a hand-designed face or audio grouping. Kenfack et al. address scarce demographic
annotations during model development~\cite{Kenfack2024}; Kallus et al. instead study disparity
assessment when the protected class is unobserved in the main data~\cite{Kallus2022}.
Putra et al. adapted ARL to continuous interview-score prediction, but their combined MSE+BCE
variant rather than MSE alone was their strongest reported First Impressions system~\cite{Putra2024}.
The present MSE comparator is therefore a disclosed architecture-controlled adaptation, not an
exact reproduction.
The key distinction in this thesis is temporal and operational: audit annotations may exist in a
sealed store, but ordinary training and proxy construction behave as if those annotations do not
exist.

These label-scarce approaches address different information problems. Kenfack et al. study
partially observed demographics and an uncertainty-aware classifier for missing sensitive
attributes~\cite{Kenfack2024}. Kallus et al. examine what can be inferred about disparity when
protected classes are unobserved in the main data but an auxiliary data source is available;
proxy-based assessment can remain only partially identified~\cite{Kallus2022}. This study does
not estimate a participant's protected identity from its clusters. Its narrower question is
whether training without access to the supplied perceived-attribute annotations changes a later,
directly joined descriptive audit. ARL supplies a distinct comparison because its high-loss
example reweighting pursues worst-case predictive performance rather than parity between two
preselected proxy score distributions~\cite{Lahoti2020}.

The proxy-fairness literature is particularly useful for stating the missing implication.
Constraining a surrogate group metric can affect an unobserved group only under assumptions
about the surrogate's relation to the group and the objective
~\cite{Gupta2018ProxyFairness}. Chen et al. assess disparity when the protected class is
unobserved~\cite{Chen2019Unawareness}. Awasthi et al. study uncertain or incomplete group
information~\cite{Awasthi2021Uncertain}, whereas Zhu et al. analyse weak proxies for missing
sensitive attributes~\cite{Zhu2023WeakProxies}. Each makes its information requirements explicit.
Robust methods for noisy protected groups also do not license substituting any convenient
cluster for an audited category~\cite{Wang2020NoisyGroups}. The present visual and audio
partitions are not imputed identities, and no theorem in these papers guarantees that their
W1 regularisation will reduce the later perceived-attribute W1.

Other demographic-free objectives target a different failure mode. Hashimoto et al. study
repeated-loss minimisation for latent high-loss groups~\cite{Hashimoto2018}; Martinez et al.
develop blind Pareto subgroup robustness~\cite{Martinez2021}. By contrast, this experiment
constrains a two-cluster
predicted-score distribution even if neither cluster has high error. Work on fairness
assessment with proxies is closer to the final audit problem, but its assumptions likewise
cannot identify this study's unlabeled clusters as actual demographic groups
~\cite{Harvey2024Proxies}. These distinctions justify evaluating both utility and a separate,
phase-aligned annotation audit rather than treating improvement of the training penalty as the
fairness result.

\section{Proxy groups and distributional regularisation}

A proxy partition can identify a region of representation space over which to impose a training
constraint. It cannot, by construction alone, identify gender, race, ethnicity or skin tone.
Face and interview audio may capture identity, the interviewer's contribution, recording
conditions, speaking style or other nuisance
variation. The thesis therefore calls the learned clusters \emph{proxy groups}, assigns neutral
identifiers, repeats construction within each training fold, and prohibits choosing a proxy by
post-hoc agreement with audit annotations.

Transfer from proxy parity to an audited category is a hypothesis, not a property of clustering.
The proxy and the perceived-attribute axis partition participants differently. Equal score
distributions across one partition can coexist with unequal distributions across the other;
reducing proxy W1 therefore gives no guarantee of reducing the audited disparity. A
sampling-only model matched to each visual or audio partition separates the effect of balanced
training batches from the additional W1 penalty. The fixed random partition tests whether an
apparent benefit depends on input-derived grouping rather than an arbitrary partition. None of
these controls identifies a cluster as a demographic class.

Wasserstein distance is useful for continuous predictions because it compares complete empirical
score distributions rather than a single threshold. Kim et al. apply a distributional objective
to continuous interview scores~\cite{Kim2023}; Yoon et al. study score-distribution alignment
for threshold-robust fair classification~\cite{Yoon2020}. For two empirical distributions with
quantile functions $F^{-1}_a$ and $F^{-1}_b$, the one-dimensional distance is
\begin{equation}
    W_1(F_a,F_b)=\int_0^1\left|F^{-1}_a(u)-F^{-1}_b(u)\right|\,du.
\end{equation}
In training, the distance is computed across a proxy partition. In the final audit, it is computed
across reportable perceived-attribute groups. These are different variables and must never be
conflated.

The one-dimensional integral is grounded in the classical CDF identity for Wasserstein
distance~\cite{Vallender1974}. Jiang et al. study score-distribution matching for
observed-group classification~\cite{Jiang2020Wasserstein}; Chiappa et al. develop a broader
optimal-transport fairness formulation~\cite{Chiappa2020OT}. For continuous outcomes, Chzhen et
al. study group-conditioned barycentres~\cite{Chzhen2020Barycenters} and then
recalibration~\cite{Chzhen2020Recalibration}.
The latter papers assume an observed, meaningful sensitive grouping. They help explain the
choice of a continuous-distribution statistic, but their statistical guarantees do not carry
over to independently learned, unlabeled visual or audio clusters.

For a binary grouping, the same one-dimensional W1 is the integrated absolute difference in
selection rates over all possible score thresholds. Kim et al. use this score-distribution
criterion for interview assessment~\cite{Kim2023}; Yoon et al. relate it to threshold-robust
classification fairness in a domain-transfer setting~\cite{Yoon2020}. It is still an
unconditional score-distribution comparison. It does not measure whether the model's errors or
residuals are similar across groups, nor does it measure an actual hiring decision. Different
fairness criteria address different aspects of a test and its use. Hutchinson and Mitchell
trace this issue through the history of test fairness~\cite{Hutchinson2019}; Verma and Rubin
compare the technical definitions that need not agree~\cite{Verma2018}. A nearly constant score
could narrow a distributional gap while
destroying useful prediction. Accordingly, the present study reports predictive utility and
groupwise error diagnostics alongside disparity, without treating lower W1 as a complete fairness
verdict. The multicategory audit also retains its categories rather than forcing them into the
binary form required by the training penalty.

Marginal parity can also conceal variation across intersections of attributes, the concern
formalised by subgroup-fairness research~\cite{Kearns2018}. This thesis does not estimate
intersectional effects: 69 independent post participants and sparse categories make that a
different, inadequately supported estimand. Reporting the available axes separately is a
transparent scope choice, not a claim that all relevant subgroups have been audited.

In personnel selection, statistical parity also has to be read beside the validity of the
selection measure and the possibility of subgroup differences in that measure
~\cite{Ployhart2008}. This thesis neither observes selection decisions nor validates job
performance. A broad survey maps bias and fairness across machine learning
~\cite{Mehrabi2021}; a more focused survey catalogues in-processing interventions
~\cite{10.1145/3551390}. Neither catalogue of
metrics resolves which construct the supplied crowd rating measures or whether a proxy partition
corresponds to an applicant group.

\section{Perceived attributes are audit constructs}

Many face datasets attach demographic-style labels to images; FairFace is a prominent example
~\cite{Karkkainen2021FairFace}. Such taxonomies are dataset-specific and visual labels cannot
establish a person's self-identified gender or ethnicity. The available audit source in this
study contains one observer's session-level perceptions of gender presentation, racialised
appearance and skin tone. The analysis preserves those source judgements, uses them only in
phase-aligned comparisons and treats ``cannot determine'' as unknown. Skin-tone annotation is
subjective and sensitive to imagery, lighting and annotator experience; evidence supporting
diverse annotators and higher replication therefore defines a limitation of this single-observer
source and an important direction for external validation~\cite{Schumann2023}.

The history of automated facial classification makes the measurement boundary consequential.
Gender Shades demonstrated intersectional accuracy disparities in commercial gender
classifiers~\cite{BuolamwiniGebru2018}. Scheuerman et al. examined how facial-analysis systems
assign gender categories~\cite{Scheuerman2019}; Khan and Fu examined how computer-vision
datasets construct racial labels~\cite{KhanFu2021}. Those studies do not measure the MIT outcome,
but they
undermine an interpretation of a visual or model-assigned category as self-reported identity.
The binary F/M entries observed in this source are a single rater's perceived presentation,
not evidence that all participants' identities are binary.

Schumann et al. studied a different, ten-category skin-tone instrument; their evidence supports
the caution about annotation conditions and replication, not the validity of this study's
project-defined five-point codes~\cite{Schumann2023}. Here the supplied labels come from one
observer, and each applies to a particular session. Their pre/post differences cannot be used as
inter-rater agreement or collapsed into a fixed participant identity. Participant-level bootstrap
intervals can quantify sampling uncertainty in the audited sample, but cannot recover the
unmeasured disagreement that independent raters might have shown. Phase-aligned joins and
sparse-category suppression keep the numerical audit faithful to what was actually annotated;
they do not repair its measurement limits.

Barrett et al. document how image context and annotator judgement affect skin-tone
annotations~\cite{Barrett2023}. Krishnapriya et al. compare manual and automated assignments
~\cite{Krishnapriya2022}, while Howard et al. examine the agreement of image-based and
self-reported skin phenotype~\cite{Howard2021}.
More recent colorimetric work offers a different instrument for improving annotation
~\cite{Cook2025}. None retrospectively validates the project-defined five-point visual
scale, and it must not be renamed as a Fitzpatrick or Monk scale. Guidance on releasing
annotator-level information reinforces the need to report the source and subjectivity of the
single available rater~\cite{Prabhakaran2021}; it does not create inter-rater evidence that
was never collected.

These labels support a bounded disparity audit of model scores. They do not support claims about
true identity, discrimination experienced by participants, legal compliance or equitable hiring.

\section{Dependent observations and small-sample evaluation}

The unit of generalisation is a participant, not an interview file. Cross-validation for
correlated data requires accounting for dependence~\cite{RabinowiczRosset2022}; repeated
measurements from one person can likewise make session-level train/test splits answer an easier,
different question~\cite{Saeb2017}. The study's participant-grouped folds prevent a
candidate's pre interview from entering training while the same candidate's post interview is
tested. These sources justify the dependence boundary; they do not independently verify the
correctness of this corpus's source-to-participant mapping.

Uncertainty must respect that same unit. Cluster-bootstrap theory treats whole clusters as the
resampling units~\cite{FieldWelsh2007}; a comparison of internal-validation bootstrap schemes
examines the consequences of clustering in risk models~\cite{Bouwmeester2013}. The
participant-cluster intervals in this
thesis quantify variation under its chosen resampling procedure. They do not cover annotation
disagreement, an alternative applicant population or unmeasured recording effects. Small-sample
cross-validation can produce large error bars~\cite{Varoquaux2018}; fold-to-fold variance is not
a general unbiased estimator of $K$-fold uncertainty~\cite{BengioGrandvalet2004}. This is why the
reported seed range, paired contrasts and intervals are interpreted together rather than
declaring a winner from one split.

Model selection on noisy validation estimates can itself bias the reported performance
~\cite{CawleyTalbot2010}. A fixed fit matrix reduces opportunities for retrospective
choice within an experiment. It does not make a follow-up independent of previously
inspected results. The original comparison, its audit-informed MIT extension and the
independent-corpus analysis therefore require distinct interpretations. These design
safeguards do not prove external validity or remedy a small participant sample.

Sparse intersectional cells can be evaluated with structured estimators, but this thesis did
not fit such a model~\cite{Herlihy2024}. Small-group disparity estimates and even the metrics
used to summarise them may be noisy or biased~\cite{Lum2022}. Simultaneous inference procedures
have been developed for multiple fairness-audit claims~\cite{CherianCandes2024}; the present
participant-percentile intervals are descriptive and are not simultaneous coverage guarantees
for every configuration-by-axis comparison. The predeclared primary contrasts therefore carry
more interpretive weight than ranking the 99 reported post-session point summaries.

\section{Research gap}

The gap addressed here is the empirical link between three questions that are often
conflated: whether a predictor works, whether a proxy constraint changes held-out scores,
and whether that change benefits independently audited groups. The original MIT
comparison isolates visual- and audio-proxy W1 from matched sampling controls. Random
proxies, adversarial reweighting and an annotation-supervised comparator provide
additional diagnostics. The linear follow-up tests the mechanism on a competitive
scorer, reports every declared penalty strength and separates training from held-out
proxy response. This contributes controlled explanatory evidence rather than a new
fairness metric or a claim that proxy parity solves hiring fairness.

The independent AVI-Personality corpus supplies a complementary interview context with
recruiter-rated outcomes and self-reported attributes~\cite{Zhang2026AVI}. That source
distinction permits external prediction and group-error analysis, but not an assumption
that an observer's appearance code and a participant's self-report measure the same
construct. A descriptive external baseline is also different from an independent
replication of the W1 intervention. The research questions preserve those boundaries
instead of pooling incompatible ratings or treating access to another corpus as
evidence of mitigation.

The provenance record is part of that contribution. Model cards call for explicit intended use
and evaluation reporting~\cite{Mitchell2019ModelCards}; datasheets document a dataset's
collection and limitations~\cite{Gebru2021Datasheets}. The sealed stages, source
mapping, annotation provenance and bounded claim ledger make these bounded empirical findings
inspectable. Documentation cannot itself validate the target or generalise the findings, but
it makes the limits of the study traceable.
